\section{Discussion}\label{sec:discussion}

The main finding is that licensing diverges after ordinary university IP-policy revisions. In the preferred documentary-reviewed sample, revisions toward more supportive communication are followed by higher licensing relative to universities that did not revise, and revisions toward more restrictive communication are followed by lower licensing. The pre-revision coefficients show no detectable difference, the gap is similar under the stricter post-support rule, and it stays positive when any single event is dropped.

How we built the sample matters for reading this result. We did not start from a small, hand-picked set of events. We started from every threshold revision, applied mechanical rules on overlap and outcome coverage, and held every eligible pair to the same documentary standards for comparability and timing. Widening the documentary review made the licensing estimate smaller than it was in the earlier event set, not larger. That does not remove selection into the direction of revision, but it makes the path from the policy corpus to the final sample much easier to follow.

The evidence is specific to one outcome. Filing and patent applications move in the same direction as licensing but much less precisely, patents issued move less, and disclosures show no corresponding increase. The licensing result also does not survive the Benjamini--Hochberg adjustment across the five outcomes. So the paper does not show that more supportive revisions improve technology transfer in general. Its strongest claim is narrower: the clearest divergence after a revision appears at the licensing stage.

That location fits the literature that treats the TTO as an organization with its own capabilities, not a passive pipe. Licensing is where office staffing, routines, the search for licensees, negotiation, and relationships with firms matter most \citep{SiegelWaldmanLink2003,ThursbyKemp2002,MarkmanEtAl2005}. It is also where continued cooperation among inventors, TTO staff, and prospective licensees matters for early-stage university technologies \citep{JensenThursby2001}. Disclosure is closer to the faculty-entry margin emphasized by \citet{JensenThursbyThursby2003}, \citet{OwenSmithPowell2001}, and \citet{BercovitzFeldman2008}. The pattern tells us where the association shows up; the data cannot tell us why it shows up there.

The companion paper outlines three ways written stance could matter \citep{SharmaPCSI2026}. Faculty may read the policy and use it to understand their obligations and the help on offer. TTO staff may carry the policy's language into forms, training, and guidance, so its stance reaches researchers who never read the document. Or the text may reflect the broader priorities of the university's leadership. If the first pathway were driving our results, disclosures should respond, and they do not. A pattern concentrated in licensing fits the second and third pathways better, although our data cannot separate them.

The content evidence is consistent with an organizational reading but plays a secondary role. Across all revisions, downward revisions are more likely to add conflict-of-interest language. That language governs faculty relationships with outside firms and could come with more involved commercialization procedures, but the presence of a keyword says nothing about legal meaning or administrative burden. Provision-level coding of the original 25 events also suggests differences in several clauses, but we did not extend those labels to the newly added events after seeing the expanded results. Testing specific textual or legal mechanisms would require independent coding of the provisions.

The main threat to identification is selection into the direction of revision, not timing. Upward and downward revisions have similar licensing trends beforehand, yet upward revisions start from much less supportive predecessor documents. The imbalance could reflect mean reversion, differences between institutions, or different reasons for revising. A university that rewrites its policy in a supportive direction may also be changing TTO leadership, budgets, workflows, or commercialization priorities. The stacked design, fixed effects, research-expenditure control, clean comparison sets, pre-trend tests, and permutation inference all improve the comparison, but none of them makes the direction of revision exogenous. The permutation test is best read as a sensitivity check that assumes exchangeable direction labels, not as randomization inference justified by the design.

It helps to keep the policy text separate from the revision episode. Institutional theory allows formal rules to drift away from practice \citep{MeyerRowan1977}, while work on administrative burden gives reasons why written procedures could change behavior \citep{MoynihanHerdHarvey2015,HerdMoynihan2018}. Our evidence rules out neither. A documented rewrite may matter in itself, or it may be the visible part of a broader organizational reform. Because we do not observe everything that changes alongside a rewrite, the post-revision contrast should be read as a feature of revision episodes, not as a structural effect of wording.

For universities and state systems, the policy lesson is modest. Ordinary IP-policy revisions deserve more attention than routine housekeeping, because revision episodes coincide with economically meaningful changes in licensing, and the direction of the communication tells us something about where those changes appear. The results do not show that rewriting a policy in more supportive language will raise licensing by itself. A stronger causal design would need variation that universities do not choose, such as a mandate from a governing system, staggered adoption of a common reform across campuses, or a legal shock that affects some institutions more than others.

The data also limit what we can conclude. The preferred sample has few upward revisions, and some event dates rest on tier-B documentary timing. AUTM reporting is voluntary, the outcomes are institution-level totals, and the survey does not follow individual inventions from disclosure to licensing. The 2023 endpoint also leaves little post-revision data for the latest events. These limits are why we report the stricter post-support sample, and why we emphasize where and when the pattern appears more than its causal size.

\section{Conclusion}\label{sec:conclusion}

This paper looks at what happens around ordinary revisions of university intellectual-property policies. We combine a text-based measure of policy communication with documentary review and AUTM technology-transfer outcomes to build an auditable set of dated revision events. We then compare upward and downward revisions with contemporaneous non-revisers in a stacked event study.

The clearest divergence after a revision is in licensing. Filing and patent applications point the same way with less precision, and invention disclosures show no matching increase. The pattern holds with stricter post-event support and when events are dropped one at a time, but it weakens once we adjust for multiple tests, and upward and downward revisions differ sharply in the stance of the policies they replace.

The paper therefore does not estimate a causal effect of policy wording. It provides evidence on ordinary policy revisions within universities, on where along the technology-transfer pipeline outcomes diverge most after a revision, and a transparent method for turning incomplete policy archives into dated events. Future work should look for reforms whose timing or direction is set from outside the university, and should code the legal and organizational changes that come with policy rewrites.

\section*{Declarations}
\noindent\textbf{Funding:} N/A\\
\textbf{Conflicts of interest:} The authors declare no conflicts of interest.\\
\textbf{Data availability:} Derived policy data, PCSI scores, revision codes, and replication code are provided in the replication package. AUTM survey data are subject to AUTM's terms of use; users may need to obtain the survey data directly from AUTM.\\
\textbf{Code availability:} The replication package reproduces the sample construction and every table and figure, except Appendix Table~\ref{tab:provisions}, which reports hand-coded provision labels; see its \texttt{README}.\\
\textbf{Author contributions:} [To be completed before submission.]
